% Vietnamese translation for OI Public Library
% Source-File: IOI中国国家候选队论文/2006/龙凡/龙凡.doc
\documentclass[11pt]{article}
\usepackage{oipl-vn}

\newcommand{\Oh}{\mathcal{O}}

\title{Nghiên cứu một lớp bài toán đoán số}
\author{Long Fan\\Trường Trung học Yali, Trường Sa}
\date{2006}

\begin{document}
\maketitle

\origtitle{Nghiên cứu một lớp bài toán đoán số}
\sourcepath{IOI China candidate team papers / 2006 / Long Fan.doc}

\renewcommand{\abstractname}{Tóm tắt}
\begin{abstract}
Bài toán đoán số là một dạng bài gần với trò chơi trong thi tin học. Một số
biến thể trả lời mỗi lần đoán bằng quan hệ lớn hơn, bằng hoặc nhỏ hơn giữa số
được đoán và số cần tìm. Bài viết nghiên cứu các dạng mở rộng của lớp bài này:
từ bài đoán số cơ bản, tới bài có giới hạn số lần trả lời theo một hướng, bài
có độ trễ phản hồi, rồi kết hợp các điều kiện ấy thành một mô hình động quy
hoạch. Qua đó, bài viết nhấn mạnh giá trị của phân tích truy hồi có hệ thống
thay vì chỉ dựa vào kinh nghiệm nhị phân.
\end{abstract}

\noindent\textbf{Từ khóa:} tìm kiếm nhị phân; truy hồi; mô hình toán học;
quy hoạch động.

\section{Bài toán xuất phát}

Bài toán cơ bản là: có một số nguyên \(X\in[1,N]\). Mỗi lần ta hỏi một số
\(Y\), người trả lời cho biết \(X<Y\), \(X=Y\), hoặc \(X>Y\). Cần tìm \(X\)
trong số lần hỏi ít nhất ở trường hợp xấu nhất.

Bản gốc xét một biến thể có hai ràng buộc:
\begin{itemize}
  \item câu trả lời của lần hỏi hiện tại chỉ đến sau khi ta đã hỏi câu tiếp
  theo, tức có độ trễ một bước;
  \item nếu ta đánh giá quá thấp số cần tìm, tương ứng nhận trả lời
  \(X>Y\), quá \(K\) lần, trò chơi kết thúc thất bại.
\end{itemize}
Với \(N\le 10^5\), \(K\le100\), cần tính chiến lược tối ưu trong trường hợp
xấu nhất.

Điểm khó là hai ràng buộc làm trực giác nhị phân không còn áp dụng trực tiếp.
Bài viết giải quyết bằng cách đi từ các mô hình đơn giản hơn và rút ra công
thức truy hồi.

\section{Bài đoán số cơ bản}

Với bài không có ràng buộc phụ, tìm kiếm nhị phân là chiến lược quen thuộc.
Nhưng để tổng quát hóa, ta nhìn theo hướng ngược: với \(i\) lần hỏi, xử lý
được đoạn dài tối đa bao nhiêu?

Gọi \(f(i)\) là độ dài lớn nhất của một đoạn có thể được xác định bằng không
quá \(i\) lần hỏi. Sau một câu hỏi, có ba trường hợp: nhỏ hơn, bằng, lớn hơn.
Hai phía nhỏ hơn và lớn hơn đều còn \(i-1\) lần hỏi; trường hợp bằng chiếm một
giá trị. Vì vậy
\[
  f(0)=0,\qquad f(i)=2f(i-1)+1.
\]
Từ đó \(f(i)=2^i-1\), và đáp án là số nhỏ nhất \(i\) sao cho
\[
  f(i)\ge N.
\]

Cách nhìn này không chỉ chứng minh lại nhị phân, mà còn chuẩn bị cho các biến
thể: thay vì đoán công thức theo kinh nghiệm, ta mô tả khả năng xử lý của
chiến lược bằng truy hồi.

\section{Giới hạn số câu trả lời theo một hướng}

Xét biến thể: nếu nhận \(K\) lần trả lời \(X>Y\), ta thất bại. Gọi
\[
  f(i,j)
\]
là độ dài lớn nhất có thể xử lý bằng \(i\) lần hỏi khi còn cho phép nhiều nhất
\(j\) lần trả lời \(X>Y\) trước khi thất bại. Nếu hỏi \(Y\), phía \(X<Y\) không
tiêu thụ hạn mức này, còn phía \(X>Y\) tiêu thụ một lần. Do đó
\[
  f(i,j)=f(i-1,j)+1+f(i-1,j-1),
\]
với các điều kiện biên
\[
  f(0,j)=0,\qquad f(i,0)=0.
\]
Đáp án là số nhỏ nhất \(i\) sao cho \(f(i,K)\ge N\).

Quan trọng hơn, công thức này cũng chỉ ra chiến lược hỏi. Nếu đang xử lý một
đoạn \([L,R]\), chọn \(Y\) sao cho bên trái có không quá \(f(i-1,j)\) phần tử
và bên phải có không quá \(f(i-1,j-1)\) phần tử. Vì hai phía không cân xứng,
đây không còn là chia đôi đơn giản.

\section{Độ trễ một bước}

Xét biến thể thứ hai: câu trả lời cho lần hỏi hiện tại chỉ nhận được sau khi
ta đã hỏi câu tiếp theo. Sau khi hỏi \(Y\), ta buộc phải hỏi thêm \(Y'\) mà
chưa biết quan hệ giữa \(X\) và \(Y\). Trong bài không có ràng buộc hướng, hai
phía là đối xứng, nên có thể giả sử \(Y'\) được đặt ở phía nhỏ hơn \(Y\).

Gọi \(g(i)\) là độ dài lớn nhất có thể xử lý bằng \(i\) lần hỏi trong mô hình
có độ trễ một bước. Một câu hỏi chính và một câu hỏi chờ phản hồi tạo ra hai
vùng có số lần hỏi còn lại khác nhau, dẫn tới truy hồi dạng Fibonacci:
\[
  g(i)=g(i-1)+g(i-2)+1,
\]
với điều kiện biên tự nhiên \(g(0)=0\), \(g(1)=1\). Vì \(g(i)\) tăng nhanh như
dãy Fibonacci, đáp án vẫn chỉ ở cấp \(\Oh(\log N)\), nhưng chiến lược hỏi
không giống tìm kiếm nhị phân thông thường.

\section{Kết hợp hai ràng buộc}

Quay lại bài toán xuất phát: vừa có độ trễ một bước, vừa giới hạn số lần nhận
trả lời \(X>Y\). Ta dùng động quy hoạch hai chiều.

Gọi
\[
  F(i,j)
\]
là độ dài lớn nhất có thể xử lý bằng \(i\) lần hỏi khi còn hạn mức \(j\) lần
trả lời \(X>Y\). Sau khi hỏi \(Y\), ta phải chọn câu hỏi chờ \(Y'\) trước khi
biết câu trả lời của \(Y\). Khác với biến thể chỉ có độ trễ, bây giờ đặt
\(Y'\) ở phía nhỏ hơn hay lớn hơn \(Y\) không còn tương đương, vì một phía tiêu
thụ hạn mức \(j\).

Do đó cần xét cả hai quyết định:
\begin{itemize}
  \item đặt câu hỏi chờ ở phía nhỏ hơn \(Y\);
  \item đặt câu hỏi chờ ở phía lớn hơn \(Y\).
\end{itemize}
Mỗi quyết định tạo ra một cách phân phối số lần hỏi còn lại và hạn mức còn
lại cho hai phía. Lấy phương án tốt hơn cho ta công thức quy hoạch động:
\[
  F(i,j)=1+\max\{\text{khả năng khi hỏi chờ bên trái},
                 \text{khả năng khi hỏi chờ bên phải}\}.
\]
Trong cài đặt, hai biểu thức bên trong được tính từ các giá trị
\(F(i-1,j)\), \(F(i-1,j-1)\), \(F(i-2,j)\), \(F(i-2,j-1)\) tương ứng với việc
một hoặc hai câu hỏi đã được dùng và việc trả lời \(X>Y\) có xảy ra hay không.

Ta tính \(F(i,j)\) theo \(i\) tăng dần, \(j=0..K\), cho tới khi
\[
  F(i,K)\ge N.
\]
Giá trị \(i\) đầu tiên thỏa điều kiện là đáp án. Từ quyết định đạt cực đại ở
mỗi trạng thái, cũng có thể khôi phục câu hỏi \(Y\) và câu hỏi chờ \(Y'\).

Với \(N\le10^5\), \(K\le100\), số trạng thái cần xét nhỏ vì \(F(i,j)\) tăng
nhanh. Bản gốc ghi nhận chương trình chạy rất nhanh trên dữ liệu trong phạm vi
đề bài.

\section{Một số bài cùng họ}

\subsection{Máy chơi đoán số bằng xu}

Trong biến thể này, nếu đoán lớn hơn \(X\) thì phải bỏ thêm \(a\) đồng xu để
tiếp tục; nếu đoán nhỏ hơn \(X\) thì phải bỏ thêm \(b\) đồng xu. Biết
\(X\in[1,N]\), cần số xu ít nhất để bảo đảm tìm được \(X\).

Gọi \(h(t)\) là độ dài lớn nhất có thể xử lý với \(t\) đồng xu. Khi hỏi một
giá trị, phía tiêu tốn \(a\) đồng còn \(h(t-a)\), phía tiêu tốn \(b\) đồng còn
\(h(t-b)\), cộng thêm trường hợp bằng:
\[
  h(t)=h(t-a)+h(t-b)+1,
\]
với \(h(t)=0\) nếu \(t<0\). Đáp án là \(t\) nhỏ nhất sao cho \(h(t)\ge N\).

\subsection{Độ trễ nhiều bước}

Một biến thể khác là câu trả lời của một câu hỏi chỉ xuất hiện sau khi đã hỏi
thêm \(c\) câu. Khi đó trạng thái không chỉ cần số lần hỏi còn lại, mà còn
phải mô tả hàng đợi các câu hỏi đang chờ trả lời. Tuy phức tạp hơn, tinh thần
vẫn giống nhau: mô tả khả năng xử lý của một trạng thái, rồi viết truy hồi
hoặc quy hoạch động từ các phản hồi có thể xảy ra.

\section{Kết luận}

Các bài đoán số thường dễ khiến ta nghĩ ngay tới nhị phân. Nhưng khi thêm các
điều kiện như giới hạn số lần trả lời theo một hướng, độ trễ phản hồi, hoặc
chi phí không đối xứng, nhị phân trực tiếp không còn đúng. Cách tiếp cận ổn
định hơn là hỏi: với tài nguyên hiện có, ta xử lý được đoạn dài nhất bao
nhiêu?

Từ câu hỏi ấy, ta xây dựng truy hồi hoặc quy hoạch động, đồng thời nhận lại
chiến lược hỏi cụ thể. Đây là giá trị chính của phương pháp: nó thay thế phỏng
đoán bằng phân tích toán học có thể mở rộng sang nhiều biến thể.

\end{document}
